\scanpage{161}
\begingroup
\makeatletter
\tagsleft@true
\makeatother
section des idéaux premiers $\unsure{\mathfrak{j}|\mathrm{j}}_y$ pour $y\in \mathrm{Y}$. Il est clair que la relation $\mathrm{Y}\subset \mathrm{Y}'$ entraîne $\mathfrak{j}(\mathrm{Y})\supset\mathfrak{j}(\mathrm{Y}')$ et que l’on a
\[
\mathfrak{j}\bigl(\bigcup_{\lambda}\mathrm{Y}_{\lambda}\bigr)=\bigcap_{\lambda}\mathfrak{j}(\mathrm{Y}_{\lambda})
\tag*{\textbf{(1.1.3.1)}}
\]
pour toute famille $(\mathrm{Y}_{\lambda})$ de parties de $\mathrm{X}$. \inserted{\handwritten{Enfin $j(\{x\})=j_x$}} \margin{\handwritten{(sert pour 1.1.7. !)}}

\textit{Proposition} \textbf{(1.1.4)}. --- (i) \textit{Pour toute partie $\mathrm{E}$ de $\mathrm{A}$, on a $\mathfrak{j}(\mathrm{V}(\mathrm{E}))=\unsure{\mathfrak{r}|\mathrm{r}}(\mathrm{E})$.}

(ii) \textit{Pour toute partie $\mathrm{Y}$ de $\mathrm{X}$, $\mathrm{V}(\mathfrak{j}(\mathrm{Y}))=\overline{\mathrm{Y}}$, adhérence de $\mathrm{Y}$ dans $\mathrm{X}$.}

(i) est conséquence immédiate des définitions et de (1.1.1.1)\,; d’autre part, $\mathrm{V}(\mathfrak{j}(\mathrm{Y}))$ est fermé et contient $\mathrm{Y}$\,; inversement, si $\mathrm{Y}\subset \mathrm{V}(\mathrm{E})$, on a $f(y)=0$ pour tout $f\in \mathrm{E}$ et tout $y\in \mathrm{Y}$, donc $\mathrm{E}\subset\mathfrak{j}(\mathrm{Y})$, $\mathrm{V}(\mathrm{E})\supset \mathrm{V}(\mathfrak{j}(\mathrm{Y}))$, ce qui prouve (ii).

\textit{Corollaire} \textbf{(1.1.5)}. --- \textit{Les parties fermées de $\mathrm{X}=\operatorname{Spec}(\mathrm{A})$ et les idéaux de $\mathrm{A}$ égaux à leurs racines (autrement dit les intersections d’idéaux premiers) se correspondent biunivoquement par les applications décroissantes $\mathrm{Y}\to\mathfrak{j}(\mathrm{Y})$, $\mathfrak{a}\to \mathrm{V}(\mathfrak{a})$\,; à la réunion $\mathrm{Y}_1\cup \mathrm{Y}_2$ de deux parties fermées correspond $\mathfrak{j}(\mathrm{Y}_1)\cap\mathfrak{j}(\mathrm{Y}_2)$, et à l’intersection d’une famille quelconque $(\mathrm{Y}_\lambda)$ de parties fermées correspond la racine de la somme des $\mathfrak{j}(\mathrm{Y}_\lambda)$.}

\textit{Corollaire} \textbf{(1.1.6)}. --- \textit{Si $\mathrm{A}$ est un anneau noethérien, $\mathrm{X}=\operatorname{Spec}(\mathrm{A})$ est un espace noethérien..}

On notera que la réciproque de ce corollaire est inexacte, comme le montre l’exemple d’un anneau intègre non noethérien ayant un seul idéal premier $\ne\{0\}$, par exemple un anneau de valuation non discrète de rang 1.

Comme exemple d’anneau $\mathrm{A}$ dont le spectre n’est pas un espace noethérien, on peut signaler l’anneau $\mathscr{C}(\mathrm{Y})$ des fonctions continues numériques sur un espace compact infini $\mathrm{Y}$\,; on sait qu’en tant qu’ensemble, $\mathrm{Y}$ s’identifie à l’ensemble des idéaux maximaux de $\mathrm{A}$, et il est facile de voir que la topologie induite sur $\mathrm{Y}$ par celle de $\mathrm{X}=\operatorname{Spec}(\mathrm{A})$ est la topologie initiale de $\mathrm{Y}$. Comme $\mathrm{Y}$ n’est pas un espace noethérien, il en est donc de même de $\mathrm{X}$.

\textit{Corollaire} \textbf{(1.1.7)}. --- \textit{Pour tout $\mathrm{x}\in \mathrm{X}$, l’adhérence de $\{\mathrm{x}\}$ est l’ensemble des $\mathrm{y}\in \mathrm{X}$ tels que $\mathfrak{j}_x\subset\mathfrak{j}_y$. Pour que $\{\mathrm{x}\}$ soit fermé, il faut et il suffit que $\mathfrak{j}_x$ soit maximal.}

\textit{Corollaire} \textbf{(1.1.8)}. --- \textit{L’espace $\mathrm{X}=\operatorname{Spec}(\mathrm{A})$ est un espace de Kolmogoroff.}

En effet, si $x$, $y$ sont deux points distincts de $\mathrm{X}$, on a, soit $\mathfrak{j}_x\not\subset\mathfrak{j}_y$, soit $\mathfrak{j}_y\not\subset\mathfrak{j}_x$, donc l’un des points $x$, $y$ n’appartient pas à l’adhérence de l’autre.

\textbf{(1.1.9)} D’après la prop. (1.1.\deleted{1}\inserted{\handwritten{2}}, (iv)), pour deux éléments $f$, $g$ de $\mathrm{A}$, on a
\[
\mathrm{D}(fg)=\mathrm{D}(f)\cap \mathrm{D}(g)
\tag*{\textbf{(1.1.9.1)}}
\]
Notons aussi que la relation $\mathrm{D}(f)=\mathrm{D}(g)$ signifie, d’après la prop. (1.1.4, (i)) et la prop. (1.1.1, (v)) que $\mathfrak{r}(f)=\mathfrak{r}(g)$, ou encore que les idéaux premiers minimaux contenant $(f)$ et $(g)$ sont les mêmes\,; en particulier, il en est ainsi lorsque $f=ug$, où $u$ est inversible.

\textit{Proposition} \textbf{(1.1.10)}. --- (i) \textit{Lorsque $\mathrm{f}$ parcourt $\mathrm{A}$, les ensembles $\mathrm{D}(\mathrm{f})$ forment une base de la topologie de $\mathrm{X}$.}

(ii) \textit{Pour tout $\mathrm{f}\in \mathrm{A}$, $\mathrm{D}(\mathrm{f})$ est quasi compact.} \inserted{\handwritten{En particulier $X=D(1)$ est \unsure{qu-cpact|qu-cpt|qu.-cpct}}}

(i) Soit $\mathrm{U}$ un ensemble ouvert dans $\mathrm{X}$\,; par définition, on a $\mathrm{U}=\mathrm{X}-\mathrm{V}(\mathrm{E})$ où $\mathrm{E}$ est une partie de $\mathrm{A}$, et $\mathrm{V}(\mathrm{E})=\bigcap_{f\in \mathrm{E}}\mathrm{V}(f)$, d’où $\mathrm{U}=\bigcup_{f\in \mathrm{E}}\mathrm{D}(f)$.
\endgroup
% archive sheet number: 160
